% \section{Contribution: Combining Knowledge Graphs and NLP to Analyze Instant Messaging Data in Criminal Investigations}
% \section{Entity Extraction from Chats Under Investigation}
\section[Entity Extraction from Investigative Chat Logs][EE from Investigative Chat Logs]{Entity Extraction from Investigative Chat Logs}
\label{sec:eechats}


%%%%%%

% During my \phd I worked with two types of documents from the Italian legal domain: civil judgements and investigative chat logs. In both cases, the main problem addressed is the difficulty that the users---such as judges or prosecutors---face in analyzing the high number of records. Indeed, my co-authored contributions seek to empower final users with tools that allow them to efficiently explore and retrieve relevant documents or chat messages.

% In the remaining of this chapter, I present the contributions of my \phd work on entity extraction in the Italian legal domain.
% Before describing these contributions, a review of the related literature on legal \ac{ie} and on \ac{ie} in Italian is provided, in order to contextualize the research.


% With respect to civil judgments, due to the lack of labeled domain data, a benchmark dataset has been annotated with a semi-automatic approach and an \ac{nel} pipeline, inspired from the one introduced in \Cref{sec:incrementalel:pipeline}, has been evaluated on the labeled data~\cite{BELLANDI2024}. Additionally, different adaptations techniques have been studied to improve \ac{ner} performance on legal data~\cite{aixia_2023}, discussing both the accuracy and the required computation.

Investigators often need to explore and extract insights from large volumes of heterogeneous documents, including \acf{ima} data. To address this need, we proposed an entity-centric approach to integrate different data sources (\ref{ob:architecture}), for enabling \acl{ia} applications such as faceted search and graph-based visualizations.

During my \phd, I collaborated with prosecutors and judicial police officers from two public prosecutor's offices in the context of two separate investigations. While the work presented in this thesis refers exclusively to the second case---an investigation into suspected corruption---\Cref{tab:chatstats} reports the statistics of the chat logs extracted and processed in both investigations, to better illustrate the volume of data that investigators must handle.

Our work was conducted under an official consultancy agreement within the framework of judicial investigations. All data, due to the private information contained, was stored locally and accessed only by authorized consultants, through a secure virtual private network (VPN) with user-password authentication; no external \acp{api} have been used at any stage.


Accordingly, this section focuses on the extraction and organization of entities from \ac{ima} data. Messages---including transcribed voice notes---were modeled within a knowledge graph designed to represent the network of contacts of the main suspect (\Cref{chap:architecture}). An improved version of the \ac{ee} pipeline introduced in \Cref{sec:eejudgements,sec:adaptation} was then applied and evaluated on a labeled corpus of chat logs. This corpus has been created through a procedure similar to the one employed for legal judgments~\Cref{sec:eejudgements}.
As in that case, the availability of labeled data for \ac{ima} analysis is limited by the presence of personal and sensitive information.

% % ok ma non qui. ne parlerò anche prima
% Our work is governed by an official investigation consultant contract. Data is stored locally with access via secure VPN and user-password authentication, and no external APIs have been used.
% Our solution processes the original data to assist investigators. However, as with any intermediary processing, whether by humans or machines, there is a risk of introducing errors, such as incomplete extractions, which could lead to bias. To mitigate this, our system is designed to allow prosecutors to verify all extracted information against the original data.

The work presented in this section was published at the \emph{25th International Conference on Web Information Systems Engineering (WISE 2024)}~\cite{pozzi_wise_2025}. Its main contributions, with respect to \acl{ee} applied to \acl{ima} data, can be summarized as follows:
\begin{itemize}
    \item Creation of an annotated benchmark dataset for investigative \ac{ima} data, enabling structured evaluation of \ac{ee} methods.
    \item Selection of a suitable speech-to-text model for transcribing voice messages.
    \item Assessment of \ac{ner} and incremental \ac{nel} pipelines on domain-specific data, identifying performance trends and key challenges.
\end{itemize}


The following sections describe the annotation process for obtaining labeled benchmark data (\Cref{sec:wise_dataset}), methodology adopted for handling voice messages (\Cref{sec:wise_multimedia}), the entity extraction pipeline (\Cref{sec:wise_entity_extraction}), and the dataset, experiments, and evaluation procedures (\Cref{sec:wise_experiments}).


% These works also include contributions in \acl{ia}, which are build upon the knowledge extracted from the domain data, and are described later in the next \Cref{chap:legalia}.


% \todo{maybe move later?}
% Data labeling is an essential part of the process required to adapt entity extraction and knowledge consolidation algorithms to the legal domain, assess their performance, and drive their improvement. In the following section, we discuss the methodology used to create an annotated benchmark from a large corpus and its features.



\subsection{Metadata Extraction}
\label{sec:wise:metdata_extraction}
Investigators usually use forensic software tools to extract the content of a seized smartphones (logical extraction)~\cite{exfiles_d1_1_2021}. In our case, we received chat logs in Microsoft Excel\footnote{\url{https://www.microsoft.com/en-us/microsoft-365/excel-c}} ``.xlsx'' format.

\Ac{ima} dumps consist of semi-structured data. In fact, unstructured message corpora are accompanied by structured metadata, which include information such as the timestamp of each message and the contact names of the sender and receiver. These metadata are highly valuable for investigators, as they allows them to reconstruct the suspect's communication network based on the number and direction of exchanged messages.  
For this reason, in this work we directly leverage chat metadata to construct a \ac{kg} that models contacts, chats, and messages, respectively through the entity classes \texttt{Person}, \texttt{Chat}, and \texttt{Message}.

\begin{figure}[ht]
\centering
\includegraphics[width=\textwidth]{images/wise/kg_model.drawio.pdf}
\caption{Knowledge Graph schema} \label{fig:model}
\end{figure}

The entire schema of the \emph{chat \acl{kg}} is depicted in \Cref{fig:model}. It incorporates all the metadata present in chat dumps and it is predisposed to host the new entities identified from an \ac{inel} pipeline---namely \texttt{Location}, \texttt{Organization}, \texttt{Date}, \texttt{Money}, and \texttt{Miscellaneous}, all of which are subclasses of \texttt{Entity}, together with \texttt{Person}. Finally, the relation \emph{mentioned in} allows persisting in which messages and chats entities have been mentioned. Consequently, investigators can trace entities to the source messages and verify the correctness of potential deduction reading the original messages (\ref{ch:traceability}). 


We consider the metadata from the dump we received to be applicable to most \acp{ima}.
For each chat they include
\begin{itemize}
  \item the list of participants with their phone numbers,
  \item the start time,
  \item and the time of the last activity.
\end{itemize}
And each message is described by
\begin{itemize}
  \item the timestamp it was sent,
  \item the name of the sender,
  \item the number of the sender,
  \item and optionally the attachments.
\end{itemize}


\begin{table}[htbp]
  \centering
  \begin{tabulary}{\textwidth}{LCCCCC}
    \hline
    \textbf{Topic} & \textbf{Chats} & \textbf{Proc. chats} & \textbf{Msgs} & \textbf{Attach. (img-audio-docs)} & \textbf{Persons} \\
    \hline
    1) Fraud & 1,133  & 801   & 45,252 & 3,324 (575-304-1,590) & 1,365 \\
    % \hline
    \textbf{2) Corruption} & 1,442 & 1,442  & 364,690 & 63,24 (51,532-6,273-4,066) & 2,351 \\
    \hline
  \end{tabulary}
  \caption{Statistics about two investigations in which I have been consultant. The work described in this section refers only to the second investigation about suspected corruption.}
  \label{tab:chatstats}
\end{table}

\Cref{tab:chatstats} shows the resulting statistics of the metadata extraction in both the investigations I was involved. The remainder of this section refers solely to the second investigate about suspected corruption.


\subsection{Benchmark Annotation}
\label{sec:wise_dataset}

\begin{table}[h]
\centering
\caption{Number of annotated mentions per entity type in the chat dataset.}
\begin{tabulary}{\textwidth}{LCCCCC}
\toprule
\textbf{Entity type} & \texttt{Person} & \texttt{Location} & \texttt{Organization} & \texttt{Date} & \texttt{Money} \\
\midrule
\textbf{Mentions} & 668 & 207 & 268 & 157 & 11 \\
\bottomrule
\end{tabulary}
\label{tab:chat_ner_stats}
\end{table}

We annotated a benchmark dataset for evaluating \ac{ner} using six chat conversations, three of which are group chats. These were selected among those sufficiently long and accompanied by audio transcripts. We adopted the same semi-automatic annotation procedure used for the judgments described in \Cref{sec:judgements_labeled_data}. Specifically, we first generated automatic annotations with the \ac{ner} component and then manually revised them using doccano~\cite{doccano}.
% In total, the dataset contains 668 mentions of \texttt{Person}, 268 of \texttt{Organization}, 207 of \texttt{Location}, 157 of \texttt{Date}, and 11 of \texttt{Money}.
The final dataset includes 1,311 annotated entity mentions, distributed across five entity types as reported in \Cref{tab:chat_ner_stats}.


\subsection{Audio Transcription} \label{sec:wise_multimedia}

% approcci confrontati: Whisper-Large, Wav2Vec, DeepSpeech
% WER 0.12 0.19 0.27
% BERTScore-F$_1$ 0.91 0.84 0.78
% transcription time of 9h audio, 1h16m 2h07m 3h02m


While we plan to extend our system to handle additional multimedia content, our experiments focused on audio, as this is the most time-consuming media type according to the investigators we collaborated with.  
To transcribe audio files, we employed Whisper (Large)~\cite{radford23a-whisper}, an \ac{asr} model developed by OpenAI. This choice followed an empirical evaluation of several \ac{asr} systems. We randomly selected 500 Italian audio files and their corresponding validated transcriptions from three sources: M-AILABS~\cite{m-ailabs_dataset}, CommonVoice~\cite{ardila-etal-2020-common}, and VoxForge~\cite{voxforge}.  

Since these datasets contain high-quality recordings, we artificially degraded the audio to better approximate the conditions of \ac{ima} voice messages, which often exhibit environmental noise and compression artifacts. We applied various distortion techniques: Gaussian noise, background noise, speed variation, pitch shifting, delay, and signal distortion.  

We evaluated the transcription quality using two complementary metrics: the \textit{Word Error Rate} (WER)~\cite{jm3} and the \textit{BERTScore}~\cite{zhang2020bertscore} (F$_1$).  
The WER measures the edit distance between the predicted transcription and the reference, normalized by the reference length~\cite{jm3}:  
\begin{equation}
  \mathrm{WER} = \frac{\text{Insertions} + \text{Substitutions} + \text{Deletions}}{\text{Number of words in Ground-truth transcript}}
\end{equation}
Lower WER values indicate better performance.  

BERTScore~\cite{zhang2020bertscore} instead measures the semantic similarity between two sentences based on contextual embeddings from a pretrained BERT model~\cite{devlin-etal-2019-bert}. It computes precision, recall, and F$_1$ by aligning words in the reference and candidate transcriptions according to cosine similarity in embedding space.
Unlike WER, BERTScore captures paraphrastic and semantic equivalence beyond exact word matching.


% approcci confrontati: Whisper-Large, Wav2Vec, DeepSpeech
% WER 0.12 0.19 0.27
% BERTScore-F$_1$ 0.91 0.84 0.78
% dati approssimati
% transcription time of 9h audio, 1h16m 2h07m 3h02m

% The results are shown in \Cref{}.

The best-performing model was Whisper (Large), which achieved a BERT score of 90.8\% and a WER of 28.2\%.



\subsection{Algorithms for Entity Extraction} \label{sec:wise_entity_extraction}

\begin{figure}
\includegraphics[width=\textwidth]{images/NEEL_chat_complete.drawio.pdf}
\caption{\Ac{inel} pipeline on chats}% \ac{ner} identifies mentions of entities,
% \underline{steve brown}, \underline{steve}, and \underline{Tom} as Person and \underline{Boston} as Location.
% \ac{nel} identifies candidate entities from the knowledge graph (KG). NIL prediction detects that \underline{steve brown}, \underline{steve}, and \underline{Tom} are linked to wrong entities and assumes they are not in the KG (NIL), while \underline{Boston} is linked to the KG. Entity clustering assigns \underline{steve brown} and \underline{steve} to the same cluster (\underline{NIL-1}), as they refer to the same entity, and \underline{Tom} to another cluster (\underline{NIL-2}).}
\label{fig:neel}
\end{figure}

Most pipeline components for entity extraction are the same as in \Cref{sec:eejudgements,sec:adaptation},
while the entity clustering has been improved for this study. 
Figure~\ref{fig:neel} depicts an example application of the \ac{inel} pipeline on chat messages.
In the remainder of this section, we describe each component of the \ac{inel} pipeline.

As the \ac{ner} component we use \emph{ITA+LGL+ICCJ900k}, the best performing model according to the adaptation study on judgments (\Cref{sec:adaptation}). It is based on the library SpaCy-transformers\footnote{https://spacy.io/universe/project/spacy-transformers} and uses the SpaCy transition-based parser with the contextualized token representations obtained from a transformer~\cite{vaswani_transformers_2017}.
% The transformer we use is based on ITALIAN-LEGAL-BERT~\cite{licari_italian-legal-bert_2022},
% % , an Italian BERT-base additionally fine-tuned with masked language modeling (MLM) for the Italian legal domain,
% additionally fine-tuned with masked language modeling (MLM) %on Italian civil judgements, and finally for \ac{ner} on civil judgements data (the model \textit{ITA+LGL+ICCJ900k}
% for the Italian legal domain and finally for \ac{ner} %TODO anon
% in~\cite{aixia_2023}. 
Using this model in the experiments can give us clues on how the best model for judgments (see \Cref{sec:eejudgements}) performs on investigative \ac{ima} data, allowing to understand the feasibility of using a single \ac{ner} model for both judgments and \ac{ima} data.


% \subsubsection{\ac{nel}}

For \ac{nel}, we use the bi-encoder architecture of \blinkita considering the \ac{kr} obtained from Italian Wikipedia---as in \Cref{sec:claw_consolidation}---additionally populated with the chat participants from the chat \ac{kg}.
To represent chat participants, we leverage BLINK zero-shot capability to encode an entity given its title and textual description. For each of them, we obtain a representation by giving the following input to the bi-encoder:
\begin{lstlisting}
[CLS] {name} [ENT] phone number: {phone1(,phone2,...)} [SEP].
\end{lstlisting}
As the majority of the people of interest for the investigation refers to entities not in Wikipedia, we link mentions of persons only to the chat participants. For organizations and locations, instead, we consider Wikipedia entities.
%are the main targets for linking in the investigative context, as the references to known organizations and locations can help disambiguate references and link messages across different chats.%\todo{controllare che sia vero}\todo{in realtà cerca di linkare anche le date}

% \subsubsection{NIL Prediction}
In the NIL prediction component, we employ the \textit{logistic regression classifier} described in \Cref{sec:claw_consolidation}.
It takes the \ac{nel} top-ranked score and \emph{secondiff} (the difference between the top-ranked score and the second-best) and produces a probability $p \in [0,1]$ as output, where 1 denotes the top-ranked link is correct and 0 assigns NIL.

        % \subsubsection{Entity Clustering and Consolidation}


The entity clustering approach, which has been improved with respect to previous work in \Cref{sec:eejudgements}, uses a supervised XGBoost classifier~\cite{chen2016xgboost}, trained with the mentions from the labeled dataset, to calculate a similarity score between pairs of mentions. The input features include lexical similarity measures, namely Jaro-Winkler~\cite{winkler1990string} and Jaccard~\cite{jaccard1901etude}, as well as semantic similarity computed as the cosine similarity of vectors from the \blinkita bi-encoder, applied to the surface forms of the two mentions. The classifier outputs a synthetic similarity score.  

Next, the pairwise similarity scores are used to create weighted graphs, where mentions are represented as nodes and the edge weights reflect the mention similarities. The final mention clusters are obtained by applying a community detection algorithm (Louvain method~\cite{blondel2008fast}) on these graphs.

%
Entity clustering is applied at document-level to group the NIL mentions referring to the same unknown entity. %For computational efficiency we do not cluster mentions across different documents.
%\todo{cosa ne facciamo delle annotazioni?}


The application of \ac{inel} has two outcomes: the chat file is annotated (\acl{sta}))
%in such a way that each annotation is associated with a span in the input text;
and the knowledge graph is updated with the discovered entities---which correspond to the entity clusters---using the ``mentioned in'' relationship (see Figure~\ref{fig:model}).
% Annotations include information about links to Wikipedia, if present, or the cluster the mention is associated with.    

% consolidation

\subsection{Experiments}
\label{sec:wise_experiments}

%The discussion of the results is based on the data acquired in the context of one investigation. 
% The results discussed in this section reflect the two-stage development of the project (see Section~\ref{sec:context}) and the maturity gap between the extraction of the enriched graph (first stage) and the application of the \ac{inel} pipeline (second stage). We consider data from the two investigations, focusing on the second one for specific analyses.  
Our evaluation has three main objectives:
\begin{enumerate}
  \item demonstrating that the combination of graph-based modeling, multimedia enrichment, and NLP-based entity extraction enhances \ac{ima} data analysis in criminal investigations through graph queries and semantic search;
  \item assessing the current quality of the proposed solution;
  \item discussing limitations and challenges.
\end{enumerate} 
      

%After a summary of statistics of the data extracted from the investigation's chat dump, we discuss the utility of the graph to get insights into investigative questions. Then, we provide more details about the extracted data and discuss the challenges of entity extraction.
%; our goal is to shed some light on the usefulness and fitness for use of the more advanced processing steps in the second phase.                 


\textbf{\ac{inel} Contribution to Knowledge Graph Enrichment.} The \ac{inel} pipeline successfully processed 1296 out of 1442 input chats, with errors caused by technical issues that can be fixed by splitting the very long chats. Statistics about the extracted entities are reported in Table~\ref{tab:stats}: the table indicates the impact of \ac{inel} on enriching the KG with the ``mentioned in'' relationship, 
%between entities and messages; 
%Statistics are reported in Table~\ref{tab:stats}. 
%Looking at the table we can demonstrate 
and the informativeness of audio transcriptions, which contain on average four times more entities than text messages, \ie 0.083 vs. 0.026.
The statistics also show that 2,361 mentions of persons have been linked to entities in the chat knowledge graph, which validates the introduction of this in-domain linking mechanism.

%In addition, it was possible to identify a very limited number of messages containing words not in common use, such as names of particular organisations or nicknames of people, which would have taken time or would have been difficult to identify as they are details that may escape a superficial first reading.

% For Person we identified 5701 entities (or clusters of mentions) of which 2055 refer to entities in the enriched knowledge graph for \ac{nel}. Note that we cluster at chat level, thus these 2055 entities may contain duplicates across different chats. With respect to Organization and Location, we respectively detected 1339 and 1892 entities, while for Date, Money, and Other we identified respectively 916, 124, and 32 mentions (here we report the number of mentions because we do not perform \ac{nel}, NIL prediction, nor clustering on these types).

% \begin{table}[h]
% % \begin{tabularx}{\textwidth}{|X|X|X|X|X|}
% \begin{tabular}{|c|c|c|c|c|c||c|c|c|c|}
% \hline
% \textbf{STATS} & \multicolumn{2}{c|}{\textbf{\#Mentions}} & \multicolumn{2}{c|}{\textbf{\#Links}} & \textbf{\#Entities} & \multicolumn{4}{c|}{\textbf{\ac{ner} Evaluation}}\\
% \hline
% & \textbf{Text}  & \textbf{Audio} & \textbf{to KB} & \textbf{NIL} && & \textbf{Prec} & \textbf{Rec} & \textbf{F$_1$} \\
% \hline
% \hline
% Per & 7765 & 520 & 2361* & 5404 & 5701 & Strong Match & 51.2 & 25.2 & 33.7 \\
% Org & 1753 & 113 & 614 & 1139 & 1339 & Strong Typed & 44.0 & 21.4 & 28.8 \\
% Loc & 2578 & 185 & 1118 & 1460 & 1892 & Partial Match & 91.9 & 46.0 & 61.3 \\
% Date & 916 & 53 & - & - & - & Partial Typed & 71.4 & 35.1 & 47.0\\
% Money & 124 & 23 & - & - & - & &  &  & \\
% Misc & 32 & 1 & - & - & - & &  &  & \\
% \hline
% % \end{tabularx}%
% \end{tabular}
% \caption{Statistics and \ac{ner} Evaluation. * Links for Person refer to the chat knowledge graph, not to Wikipedia.}
% \label{tab:stats eval}
% \end{table}

\begin{table}[htbp]
    \centering
    % \begin{tabular}{|c|c|c|c|c|c|}
    \begin{tabular}{l|cc|cc|c}
      \toprule
      \textbf{Type} & \multicolumn{2}{c|}{\textbf{Mentions}} & \multicolumn{2}{c|}{\textbf{Links}} & \textbf{Entities} \\
      & \textbf{Text}  & \textbf{Audio} & \textbf{to \ac{kr}} & \textbf{NIL} & \\
      \midrule
      % \hline
      \texttt{Person} & 7765 & 520 & 2361* & 5404 & 5701 \\
      \texttt{Location} & 2578 & 185 & 1118 & 1460 & 1892 \\
      \texttt{Organization} & 1753 & 113 & 614 & 1139 & 1339 \\
      \texttt{Date} & 916 & 53 & -- & -- & -- \\
      \texttt{Money} & 124 & 23 & -- & -- & -- \\
      \texttt{Miscellaneous} & 32 & 1 & -- & -- & -- \\
      \bottomrule
    \end{tabular}
    \caption{Statistics of the \ac{inel} extraction. *Links for Person refer to the chat knowledge graph, not to Wikipedia.}
    \label{tab:stats}
\end{table}

\begin{table}[htbp]
    \centering
    \begin{tabular}{lccc}
      \toprule
      & \textbf{Precision} & \textbf{Recall} & \textbf{F$_1$} \\
      \midrule
      Strong-typed & 44.0 & 21.4 & 28.8 \\
      Partial-typed & 71.4 & 35.1 & 47.0 \\
      \bottomrule
    \end{tabular}
    \caption{\ac{ner} Evaluation. \textit{Strong} counts perfect matches of span and type. \textit{Partial} counts as correct when there is an overlap between predictions and ground truth.}
    \label{tab:eval}
\end{table}


% audio menzioni
% {'persona': 520,
%  'luogo': 185,
%  'organizzazione': 113,
%  'denaro': 23,
%  'data': 53,
%  'altro': 1}

% ner stats
% {'data': 916,
%  'organizzazione': 1753,
%  'persona': 7765,
%  'luogo': 2578,
%  'denaro': 124,
%  'altro': 32}

% linking_stats

% {'neo4j': 2343, 'organizzazione': 614, 'persona': 2361, 'luogo': 1118}

% nil_stats

% {'organizzazione': 1139, 'persona': 5404, 'luogo': 1460}

% Selection deleted
% cluster_stats

% 'neo4j': 2055,
% {'data': 793,
%  'organizzazione': 1339,
%  'persona': 5701,
%  'luogo': 1892,
%  'denaro': 100,
%  'altro': 28}

% audio menzioni
% {'persona': 520,
%  'luogo': 185,
%  'organizzazione': 113,
%  'denaro': 23,
%  'data': 53,
%  'altro': 1}

% performance P R F$_1$
% strong 0.511737089	0.251538462	0.337287261
% strong typed 0.439749609	0.216153846	0.289840124
% approx 0.683881064	0.33954934	0.453790239
% approx typed 0.561815336	0.276792598	0.370867769
% partial 0.918622848	0.464031621	0.616596639
% partial typed 0.713615023	0.353488372	0.472783826


%%% old investigation

% In these chats, the \ac{inel} pipeline found \todo{update}7486 mentions of entities: 4592 of type ``person", 1166 of type ``location", 606 of type ``date", 229 of type ``money" and 859 of type ``organization". This means that, thanks to NLP, we can enrich the knowledge graph with 7486 ``mentioned in" relationships between ``entity" nodes and the ``message" nodes. We should note that the addition of the transcriptions of voice-messages 

% that contain actual messages was used to quantify the number of the ``mentioned in" relationships that could be added to the knowledge graph. In these chats, the \ac{inel} pipeline found \todo{update}7486 mentions of entities: 4592 of type ``person", 1166 of type ``location", 606 of type ``date", 229 of type ``money" and 859 of type ``organization". This means that, thanks to NLP, we can enrich the knowledge graph with 7486 ``mentioned in" relationships between ``entity" nodes and the ``message" nodes.
%%%% old investigation



% We should note that the addition of the transcriptions of voice-messages 
%to the documents 
% helped to recognize 138 more entities than when the pipeline was applied to the plain-text messages alone. This suggests that voice messages may have a higher information content, since the predicted average number of ``entity" nodes connected to them by a ``mentioned in" relation (0.24) is higher than the average number of ``entity" nodes connected to a ``Text-message" node (0.16).
%At the time of writing this paper, we are still working on finalizing the interconnection between the \ac{inel} pipeline and the knowledge graph. For this reason, we can't yet provide any additional statistics on the number of different ``entity" nodes.\todo{Per Prof: controllare se OK}
% An example of how the \ac{inel} pipeline identifies entities in a transcription is shown in Figure~\ref{fig:\ac{ner} example}. 
%Please note that this 
% The screenshot is taken directly from DAVE and shows how the annotated document is presented to the user. In this example, several mentions of locations and organizations were found in a voice message that referred to a meeting (``incontro" - one of the keywords searched to answer the investigative questions). 
%, in a chat that had references to a "meeting" (a relevant topic in ou a voice message related  

% \begin{figure}[h]
% \centering
% \includegraphics[width=0.8\textwidth]{\ac{ner} example.png}
% \caption{Example entity annotations on a transcription. ``Persona" refers to entity of type ``Person", ``Luogo" to ``Place", ``Organization" to ``Organization" and ``Data" to ``Date"} \label{fig:\ac{ner} example}
% \end{figure}

\textbf{Preliminary Insights on the Quality of \ac{inel} Annotations.}
We discuss a first evaluation of the \ac{ner} component. 
%, the task for which a comparison with our previous work is easier.  
%To quantify the quality of the NLP pipeline 
% We annotated a labeled dataset for evaluating \ac{ner} from six chats of which three are group chats, selected from the ones long enough and with audio transcripts. %We used the same methodology used in previous work~\cite{BELLANDI2024105904,aixia_2023}: % TODO anon riaggiungere citazione
% We computed a first set of automatic annotations using the \ac{ner} component and thoroughly revised them with doccano~\cite{doccano}. At the end of the process, we annotated 668 mentions of Person, 268 of Organization, 207 of Location, 157 Dates, and 11 Money.

% 78 mentions of time, 13 phone numbers, 18 website addresses, 2 email addresses, and 12 misc. However, the \ac{ner} evaluation only considers Person, Organization, Location, Date, and Money since they are the types currently extracted by the \ac{ner} module.

% gs stats
% {'persona': 668,
%  'luogo': 207,
%  'organizzazione': 268,
%  'data': 157,
%  'numerocell': 13,
%  'altro': 12,
%  'sito': 18,
%  'denaro': 11,
%  'email': 2,
%  'time': 78}

% of  for four of the longest chats and 
% %. We proceeded to evaluate the quality of the Named Entity Recognition algorithm, by 
% %In particular, 
% %With respect to precision, we found some mistakes on transcriptions related to I'm  

% Results shown in Table~\ref{tab:eval} are significantly worse than those obtained on other domain-specific data on which the pipeline we adapted was tested~\cite{ijckg2022,aixia_2023}. This confirms that \ac{inel} on \ac{ima} data is challenging because of the specific data distributions and suggests that in-distribution fine-tuning is necessary. % to detect entities. % TODO anonymized. valutare se reinserire
Results shown in Table~\ref{tab:eval} are not satisfactory, confirming that \ac{inel} on \ac{ima} data is challenging because of the specific data distributions and suggesting that in-distribution fine-tuning is necessary.
Metadata are precisely identified enabling accurate filtering with faceted search.
% Indeed \ac{ima} data contain peculiarities such as person name abbreviations or inconsistent capitalization of names and locations. In relation to fitness for use, we make two remarks. 1) Metadata are precisely identified enabling accurate filtering in DAVE, for instance, by message sender. 2) DAVE supports text search, which makes it possible to complement entity-based retrieval to improve result recall.        


%\begin{itemize}
%    \item Questions to be answered with the report, vs. type of queries vs example of queries
%    \item How these queries are empowered by NLP (exemplify) and speech-to-text (keyword search like "incontro") 
%    \item Statistics on a chat sample: how many entities are found? How many same-as links in the graph? How many speech-to-text transcriptions? 
%    Statistics that are relevant: mentions divided by class; how many are linked to Wikipedia, how many clusters, how many entities found in text are associated with entities in the KG, and how many are novel? Problem: one chat as an example, or more chats with average.  \\
%\item Some insights on the performance of NLP on chats. \ac{ner}: we count one if the correct annotation is included or includes our annotation AND class is correct (approximate typed mention match). Linking: we count how many entities for which a link is found by the algorithms have a correct link. Most recent measures proposed to evaluate clustering are quite complex to compute in post-hoc conditions with a manual approach, where we need to pick up one cluster at a time and control its quality; to give some insights on the quality of clusters, we use a heuristic measure defined as follows: we sample x clusters, each of which is associated with an entity label; we count how many mentions in the cluster are actually referred to this target entity. Observe that, in this way, we only evaluate if mentions in a cluster are correctly associated, while we do not consider the capability of joining together all mentions referred to the same entity to one cluster; however, in this context of semantic search, we believe that a measure inspired by precision is more relevant than a measure inspired by recall (in fact a user can also explore more cluster with a same label).      
%\end{itemize}


%\todo{forse meglio mettere una sezione tipo challenges. e citare qualche challenge anche nell'intro.}

        

% \textbf{Limitations and Challenges.}
% Results in Table~\ref{tab:eval} confirm that \ac{inel} on \ac{ima} data is challenging due to the peculiar distribution of entities, indicating that in-distribution fine-tuning is required to attain competitive performance. At the same time, metadata are reliably identified, enabling accurate filtering via faceted search. To overcome the weaknesses of pretrained models in the \ac{inel} pipeline, the labeled benchmark must be enlarged and the models fine-tuned on \ac{ima} data.%; based on our previous experience~\cite{ijckg2022}, we expect a significant improvement across all performance measures. % TODO anon. valutare se riaggiungere
%
%The \ac{inel} stack should be better customized for the domain, \eg via in-domain fine-tuning. %Clustering is attracting the attention of the research community and we believe that investigations and \ac{ima} data provide a representative and challenging domain to develop better algorithms for this task. 
%\todo{forse potremmo togliere questa discussione sul clustering visto che non è tanto esplorato nel paper}
%In the customization of entity clustering in the domain of \ac{ima}, we also need to improve the interconnection between the knowledge graph and the \ac{inel} pipeline.
%For this reason, we can't yet provide any additional statistics on the number of different ``entity" nodes.\todo{Per Prof: controllare se OK}
% The integration of generative LLMs is still limited.
% The QA module requires significant improvement and
%We did not include image-to-text enrichment yet, a highly desired features.
